\begin{definitionbox}{Definition --- Metric}
\begin{definition}[Metric]\label{def:metric-on-set}
Let \(X\) be a set.  A \emph{metric} on \(X\) is a real-valued function
\[
  d : X \times X \to \mathbb{R}
\]
such that for all \(x,y,z\in X\):
\begin{enumerate}
  \item \(d(x,y)\geq 0\), and \(d(x,y)=0\) if and only if \(x=y\);
  \item \(d(x,y)=d(y,x)\);
  \item \(d(x,z)\leq d(x,y)+d(y,z)\).
\end{enumerate}
\end{definition}
\end{definitionbox}

\begin{remark*}[Standard quantified statement]
\[
\begin{aligned}
&d:X\times X\to\mathbb{R}
\\
&{}\land
\forall x,y,z\in X{:}\;
d(x,y)\geq 0
\land
\bigl(d(x,y)=0\Longleftrightarrow x=y\bigr)
\\
&{}\land
d(x,y)=d(y,x)
\land
d(x,z)\leq d(x,y)+d(y,z).
\end{aligned}
\]
\end{remark*}

\begin{remark*}[Predicate reading]
\[
d(x,y)
\]
means the distance assigned by \(d\) from \(x\) to \(y\).
\end{remark*}

\begin{remark*}[Interpretation]
A metric assigns distances between points of a set.  Its axioms say that
distance is nonnegative, that zero distance identifies exactly equal points,
that distance is symmetric, and that the direct distance from \(x\) to \(z\)
is no larger than the distance obtained by passing through \(y\).
\end{remark*}

\begin{remark*}[Historical note]
This definition is modeled on Definition~1.1.1 of
{\'O} Searc{\'o}id's \emph{Metric Spaces} \cite{OSearcoidMetricSpaces}.
\end{remark*}

\LeanFormalizes{def:metric-on-set}{lra-lean}{LRA.VolumeIV.MetricSpaces.Foundations.MetricSpace}{MetricDefinition}{statement}

\DefinitionalRoot

\begin{definitionbox}{Definition --- Metric Space}
\begin{definition}[Metric Space]\label{def:metric-space}
A \emph{metric space} is a pair \((X,d)\), where \(X\) is a set and \(d\) is
a metric on \(X\).
\end{definition}
\end{definitionbox}

\begin{remark*}[Standard quantified statement]
\[
(X,d)
\quad\text{with}\quad
d \text{ a metric on } X.
\]
\end{remark*}

\begin{remark*}[Interpretation]
A metric space is a set equipped with a chosen metric.  The set supplies the
points, and the metric supplies the distance structure on those points.
\end{remark*}

\begin{remark*}[Historical note]
This definition is modeled on Definition~1.1.1 of
{\'O} Searc{\'o}id's \emph{Metric Spaces} \cite{OSearcoidMetricSpaces}.
\end{remark*}

\LeanFormalizes{def:metric-space}{lra-lean}{LRA.VolumeIV.MetricSpaces.Foundations.MetricSpace}{MetricSpaceDefinition}{statement}

\begin{dependencies}
\begin{itemize}
  \item \hyperref[def:metric-on-set]{Definition: Metric}
\end{itemize}
\end{dependencies}

\begin{theorembox}{Theorem}
\begin{theorem}[Rearrangement of the Triangle Inequality]
\label{thm:metric-rearranged-triangle-inequality}
Let \((X,d)\) be a metric space.  For all \(a,b,c\in X\),
\[
  \bigl|d(a,b)-d(b,c)\bigr|\leq d(a,c).
\]

\smallskip
\noindent
\hyperref[prf:metric-rearranged-triangle-inequality]{\textit{Go to proof.}}
\end{theorem}
\end{theorembox}

\begin{remark*}[Standard quantified statement]
\[
\forall a,b,c\in X{:}\;
\bigl|d(a,b)-d(b,c)\bigr|\leq d(a,c).
\]
\end{remark*}

\begin{remark*}[Predicate reading]
\[
\bigl|d(a,b)-d(b,c)\bigr|\leq d(a,c)
\]
means that the two distances from \(b\) to \(a\) and from \(b\) to \(c\)
differ by at most the distance from \(a\) to \(c\).
\end{remark*}

\begin{remark*}[Interpretation]
In a metric space, the distances from \(b\) to \(a\) and from \(b\) to \(c\)
cannot differ by more than the distance between \(a\) and \(c\).  This is the
triangle inequality rearranged so that one distance controls the difference
between two others.
\end{remark*}

\begin{remark*}[Historical note]
This theorem is modeled on Theorem~1.1.2 of
{\'O} Searc{\'o}id's \emph{Metric Spaces} \cite{OSearcoidMetricSpaces}.
\end{remark*}

\LeanFormalizes{thm:metric-rearranged-triangle-inequality}{lra-lean}{LRA.VolumeIV.MetricSpaces.Foundations.MetricSpace}{rearrangement_of_triangle_inequality}{statement}

\begin{dependencies}
\begin{itemize}
  \item \hyperref[def:metric-space]{Definition: Metric Space}
\end{itemize}
\end{dependencies}
